\section{Radiation pressure}
\label{sec:radiation-pressure}

This section documents the radiation-pressure feedback channel
(\code{GEARFeedback:radiation\_pressure\_efficiency},
\code{radiation\_policy\_radiation\_pressure}), a separate mechanism
from the HII photoionization budget documented in the algorithm
section of the HII photoionization theory document
(\code{radiation.tex}/\code{radiation.pdf}, this same directory):
photoionization tags and re-floors gas
temperature, while radiation pressure injects momentum directly into
gas neighbouring the star. The two channels draw on \emph{different}
integrals of the star's spectrum ($Q_H$ vs.\ $L_{\rm bol}$, see
Section~\ref{sec:rp-vs-hii} below) and use the same SPH gather around
the star, but are otherwise independent: a star can have one policy
enabled without the other. Not yet covered anywhere else in this
document; written up here for the first time (2026-07-19).

\subsection{Physical distinction from photoionization}
\label{sec:rp-vs-hii}

Two genuinely different physical mechanisms in a real star, not two
bookkeeping categories of the same effect (clarified 2026-07-20).

\textbf{Photoionization} is a quantum/atomic process: a photon with
energy above $13.6\,$eV (Lyman continuum, $\lambda<912\,$\AA) is
absorbed by a bound electron in hydrogen, ejecting it; the photon's
energy above the ionization threshold becomes the photoelectron's
kinetic energy, which thermalizes via collisions; this is what heats
the gas to $\sim10^4\,$K, balanced against recombination/line cooling.
Only photons above that specific energy threshold count, and the
relevant budget is $Q_H$ (ionizing photon \emph{number} rate, the
blackbody-fit section of the HII photoionization theory document), not energy.
This is fundamentally a \textbf{thermal/chemical} effect: it sets the
gas's ionization state and temperature, and in this code drives
dynamics only indirectly, through the resulting thermal-pressure
contrast against the cold ambient medium (the classical Spitzer D-type
expansion that STARBENCH, validated in the HII photoionization
logbook's own validation section, checks).

\textbf{Radiation pressure} is a mechanical process: \emph{any}
absorbed or scattered photon, at any wavelength, transfers momentum
$p=E/c$ to whatever absorbs or scatters it. This draws on $L_{\rm
bol}$: the star's entire spectrum, dominated for most stars by the
non-ionizing UV/optical continuum, not the ionizing band specifically
(the mass-luminosity fit vs.\ the separate ionizing-rate fit behind
$Q_H$, both in the HII photoionization theory document's blackbody-fit
section. Two different fits of the same underlying blackbody, not the same
fit reused). It is
a \textbf{dynamical} effect: a direct push, independent of whether
the interaction happens to ionize anything, and in this code it acts
through Eq.~\ref{eq:rp-momentum-rate} directly, with no thermal
intermediary.

Two consequences worth keeping precise:
\begin{itemize}
  \item \textbf{They are not disjoint at the photon level.} An
    ionizing-photon absorption \emph{also} transfers momentum
    $h\nu/c$ in that same event: momentum conservation does not care
    about the atomic-physics outcome. But that ionizing-photon
    momentum is a small slice of the total budget, since $Q_H$
    photons are only part of $L_{\rm bol}$ for most stars, and dust
    reprocessing/IR trapping (the $(1+\tau_{\rm IR})$ term,
    Eq.~\ref{eq:rp-tauIR}) adds momentum from photons that were never
    ionizing to begin with. This code does not attempt to avoid
    double-counting this overlap; see the open interaction question
    in the GEAR radiation logbook's radiation-pressure open-questions
    section.
  \item \textbf{They drive HII region expansion through different,
    competing channels.} Photoionization heats gas against a cold
    ambient medium ($\sim10$--$100\,$K): the resulting
    \emph{thermal-pressure} contrast drives the D-type expansion.
    Radiation pressure drives expansion through \emph{direct momentum}
    deposition instead, no heating required.
    \citet{KrumholzMatzner2009}, already cited above for the Tier-2
    validation proposal, is literally a paper about comparing these
    two channels: thermal-pressure-driven expansion dominates for the
    modest OB associations STARBENCH-style tests represent, while
    radiation pressure only becomes the dominant \emph{dynamical}
    driver for far more luminous, compact sources (starburst clusters,
    super star clusters) where $L$ is large and the region is
    dense/compact enough for $\tau_{\rm IR}$ to matter.
\end{itemize}

\subsection{Momentum injection: single-scattering plus IR trapping}
\label{sec:rp-momentum}

Each active star's momentum injection rate is
(\code{radiation\_get\_star\_physical\_radiation\_pressure},
\code{src/feedback/GEAR/radiation.c}):
\begin{equation}
  \dot p_{\rm rad} = \frac{L_{\rm bol}}{c}\,(1+\tau_{\rm IR}),
  \label{eq:rp-momentum-rate}
\end{equation}
the standard single-UV-scattering-plus-dust-IR-trapping recipe: direct
absorption/scattering of the star's own UV/optical photons deposits
$L_{\rm bol}/c$ of momentum per unit time; photons re-radiated in the
IR by dust can scatter multiple times in an optically thick shell
before escaping, boosting the momentum transfer by a further factor of
$(1+\tau_{\rm IR})$ over a single scattering.

\textbf{Citation confirmed (2026-07-23), direct primary-source check.}
Darwin confirmed this codebase's implementation is transcribed from
\citet{Hopkins2020} (the FIRE-2 radiative-feedback methods paper,
arXiv:1811.12462, distinct from the FIRE-2 numerics paper
\citet{Hopkins2018} already cited elsewhere in this document). Fetched
and grepped the full text directly: \citet{Hopkins2020}'s Appendix~A
gives \textbf{exactly}
$(\kappa_{\rm FUV},\kappa_{\rm NUV},\kappa_{\rm opt},\kappa_{\rm IR}) =
(2000, 1800, 180, 10)\,(Z/Z_\odot)\,{\rm cm^2\,g^{-1}}$: this
codebase's $\kappa_{\rm IR}=10\,{\rm cm^2\,g^{-1}}\,(Z/Z_\odot)$
(Eq.~\ref{eq:rp-tauIR}) matches precisely, not just order-of-magnitude.
\citet{Agertz2013}/\citet{Marinacci2019} remain plausible independent
lineages sharing the same functional form but are no longer the
open question: \citet{Hopkins2020} is confirmed as the actual source.

\textbf{New finding from the same primary-source check: a term this
codebase is missing.} \citet{Hopkins2020} \S4.4, footnote 16 gives the
full LEBRON momentum-coupling formula as
$\dot p = \tilde\tau\,L_{\rm bol}/c$, with
$\tilde\tau = (1-e^{-\tau_{\rm single}})\,(1+\tau_{\rm eff,IR})$,
\emph{not} simply $(1+\tau_{\rm IR})$ as
Eq.~\ref{eq:rp-momentum-rate} implements. The
$(1-e^{-\tau_{\rm single}})$ factor is the fraction of the
\emph{input} (non-ionizing UV/optical) luminosity actually absorbed
before reaching the dust that reprocesses it into the trapped IR
term: this codebase's formula implicitly assumes
$\tau_{\rm single}\to\infty$ (i.e.\ $(1-e^{-\tau_{\rm single}})=1$,
full absorption), which only holds deep in a dense cloud. Hopkins et
al.\ report their own galaxy-scale simulations couple only
$\langle\tilde\tau\rangle\sim0.4$--$0.7$ of the bolometric luminosity
on average (an order-unity fraction of sightlines are optically thin
even in the optical); so at Darwin's production resolution
(10\,$M_\odot$ to a few $10^4\,M_\odot$ per gas particle), where much
of the coupling volume around a star can be diffuse, this
missing factor likely causes this codebase to \emph{over}-estimate
direct/single-scattering radiation pressure, on top of (and largely
independent from) the $\kappa_{\rm IR}$ bug's IR-trapping
\emph{under}-estimate (GEAR radiation logbook, ``Bugs found and
fixed (2026-07-23)'' section): the two
are separate errors, not a cancelling pair, and there is no reason to
expect them to offset in general. Implementing this requires an
opacity for the input band(s), which \citet{Hopkins2020}'s own
$\kappa_{\rm FUV}/\kappa_{\rm NUV}/\kappa_{\rm opt}$ spread by over an
order of magnitude (2000 vs.\ 180\,${\rm cm^2\,g^{-1}}$): a single
effective $\kappa_{\rm single}$ (e.g.\ luminosity/temperature-weighted
via the blackbody fit already computed for $T_K$, in the HII
photoionization theory document) would capture the leading effect without adopting
\citet{Hopkins2020}'s full 5-band architecture. Not yet implemented;
this is judged the single most valuable remaining physics gap in this
channel, ahead of the resolution/banding items below.

$L_{\rm bol}$ is \code{sp->feedback\_data.radiation.L\_bol}, set from
the same blackbody mass--luminosity fit as the ionizing photon rate
(the HII photoionization theory document's blackbody-fit section) and scaled
by \code{radiation\_pressure\_efficiency} (\code{feedback\_common.c}):
an overall dial on this channel independent of the photoionization
budget, default behaviour unaffected unless explicitly set $\ne1$.

The IR optical depth is
\begin{equation}
  \tau_{\rm IR} = \kappa_{\rm IR}\,\Sigma_{\rm gas}, \qquad
  \kappa_{\rm IR} = 10\,\frac{\rm cm^2}{\rm g}\,\frac{Z}{Z_\odot},
  \label{eq:rp-tauIR}
\end{equation}
(\code{radiation\_get\_physical\_IR\_optical\_depth},
\code{radiation\_get\_physical\_IR\_opacity}): a metallicity-linear
dust opacity of $10\,{\rm cm^2\,g^{-1}}$ at solar metallicity, a
standard order-of-magnitude value for dust-processed IR opacity in this
class of model. $Z$ here is \code{sp->feedback\_data.Z\_star}, the
SPH-kernel-weighted gas metallicity at the star's own location
(accumulated in \code{radiation\_iact\_nonsym\_feedback\_density}
alongside $\nabla\rho_\star$ below, normalized by the same
$1/h^{d+1}$ factor in \code{feedback\_prepare\_radiation\_feedback}
that finishes $\nabla\rho_\star$; both are zero, not divided-by-zero,
for a star with no gas neighbours this step, since
\code{enrichment\_weight} being zero is checked before the divide).

\subsection{Gas column density: Sobolev approximation}
\label{sec:rp-sobolev}

$\Sigma_{\rm gas}$ in Eq.~\ref{eq:rp-tauIR} uses the Sobolev
approximation, the same length-scale trick used for line-of-sight
column densities in stellar-atmosphere and ISM radiative-transfer
codes:
\begin{equation}
  \Sigma_{\rm gas} = \rho_\star \times
  \left(h_\star\,\gamma_k + \ell_{\rm Sobolev}\right), \qquad
  \ell_{\rm Sobolev} = \frac{\rho_\star}{|\nabla\rho_\star|},
  \label{eq:rp-sobolev}
\end{equation}
(\code{radiation\_get\_comoving\_gas\_column\_density\_at\_star}), i.e.\
the local gas density times an effective path length: the kernel
support radius $h_\star\gamma_k$ (a floor: always included, not just
a fallback) plus the Sobolev length $\rho/|\nabla\rho|$, the
characteristic distance over which the density itself changes by a
factor of $e$. $\rho_\star$ is \code{enrichment\_weight} (the star's
own SPH-averaged local gas density, already computed by the regular
enrichment density loop, reused here rather than a separate
accumulator); $\nabla\rho_\star$ is
\code{sp->feedback\_data.grad\_rho\_star}, an ordinary SPH kernel-gradient
density estimate accumulated in the same density-loop pass as $Z_\star$
above. \textbf{Degenerate case handled:} a locally uniform density
field (zero gradient, e.g.\ unperturbed glass/grid initial
conditions, or a star with too few neighbours to resolve a gradient)
makes $\ell_{\rm Sobolev}$ formally undefined (0/0); the code checks
$|\nabla\rho_\star|>0$ explicitly and falls back to the kernel-radius
term alone in that case, rather than propagating a NaN.

\textbf{Confirmed against the primary source (2026-07-23).}
\citet{Hopkins2020} \S4.4 gives their own local column-density
estimate as $\langle\Sigma\rangle \approx \rho\,[h + \rho/|\nabla\rho|]$,
with $h=(m_i/\rho)^{1/3}$ the local resolution-element size,
structurally identical to Eq.~\ref{eq:rp-sobolev}'s
$\rho_\star(h_\star\gamma_k + \ell_{\rm Sobolev})$: same
floor-plus-Sobolev-length sum, just a slightly different (arguably
more physically motivated, since it is the actual SPH kernel support
radius rather than a generic mass/density cube root) floor term.
\textbf{One correction to an earlier assumption in this discussion}:
Darwin recalled FIRE capping the $\rho/|\nabla\rho|$ term at the
smoothing length; the formula as quoted in \citet{Hopkins2020} \S4.4
is a plain sum, no $\min(\cdot)$: i.e.\ no cap is visible in this
specific paper. This may be in \citet{Hopkins2018} (their ``Paper~I'',
not yet checked) or a later code update, or simply a misremembering;
not resolved this pass. Practically: this codebase's formula already
matches \citet{Hopkins2020}'s own documented approach as-is, so the
resolution-robustness concern raised below (production runs at
10--$10^4\,M_\odot$ per particle) is about noise in $|\nabla\rho|$
at coarse resolution, not about deviating from the primary source:
a cap would be a robustness \emph{improvement over} what
\citet{Hopkins2020} documents, not a correction to match it.

\subsection{Distributing the kick to gas particles}
\label{sec:rp-distribute}

$\dot p_{\rm rad}$ (Eq.~\ref{eq:rp-momentum-rate}) is a single
star-level rate; distributing it to individual gas neighbours uses the
same SPH-kernel-weighted-share pattern already used for chemical
enrichment (\code{radiation\_iact\_nonsym\_feedback\_apply}):
\begin{equation}
  \Delta\vec p_j = -\,\dot p_{\rm rad}\,\Delta t_\star\,
  \frac{m_j w_{ij}}{\rho_\star}\,\hat{\vec r}_{j\star},
  \label{eq:rp-per-particle}
\end{equation}
where $\Delta t_\star$ is the star's own current timestep, $w_{ij}$ the
kernel weight, $m_j$ the gas particle's mass, and $\hat{\vec r}_{j\star}$
the unit vector pointing radially outward from the star to the gas
particle (the code computes $-dx/r$ with $dx \equiv \vec x_\star -
\vec x_j$, i.e.\ away from the star, with the usual comoving/physical
$\times\,a$ conversion applied). The result accumulates into
\code{xpj->feedback\_data.radiation.delta\_p}: a field distinct from
the top-level \code{xpj->feedback\_data.delta\_p} used for
supernova/wind momentum (\code{feedback\_struct.h}), applied separately
in \code{feedback\_update\_part\_radiation}
(\code{feedback\_data\_iact.h}: see the bug fixed in the HII
photoionization logbook's bugs-and-fixes section,
\code{hit\_by\_radiation} was never actually
set until 2026-07-19, so this accumulated momentum was silently
discarded every step until then).

